% Cuerpo científico recuperado y distribuido en su residencia terminal.
% Cuerpo editor-ready, sin preámbulo y sin portada autónoma.
% Residencia propuesta: inmediatamente después del teorema del doble círculo
% y antes de Navier--Stokes. No se consume como sustituto de RH ni del
% desarrollo íntegro del doble círculo.

\section{Entrelazamiento exacto de la holonomía nonádica y la fase espectral}
\label{sec:cpe-entrelazamiento-exacto}

Sea \(\mathcal H_W\) el espacio de Hilbert obtenido de la forma positiva de
Guinand--Weil, sea \(H_W\) el generador autoadjunto del grupo de
traslaciones y sea
\[
 \mathcal C_W=(H_W+i/2)(H_W-i/2)^{-1}
\]
su transformada de Cayley. Para cada profundidad \(k\), sea \(Z_k\) el
conjunto de estados cilíndricos enriquecidos, sea \(\mu_k\) su medida
proyectiva de soporte completo y sea
\(\tau_{k+1,k}:Z_{k+1}\to Z_k\) el truncamiento. Se considera
\[
 \mathscr K_k=\ell^2(Z_k)\widehat\otimes\mathcal H_W,
 \qquad
 \Omega_k=\sum_{z\in Z_k}\sqrt{\mu_k(z)}\,e_z,
 \qquad
 J_k\psi=\Omega_k\otimes\psi.
\]
Si \(z'\) es descendiente de \(z\), se toma
\(w_k(z'\mid z)=\sqrt{\mu_{k+1}(z')/\mu_k(z)}\). Para una ruta
\(\gamma:z\to z'\), el cociclo entero de memoria se denota por
\(m(\gamma)\), y el transporte torcido se define por
\begin{equation}
 \widetilde U_k(e_z\otimes\psi)
 :=\sum_{\tau_{k+1,k}z'=z}
 w_k(z'\mid z)\,e_{z'}\otimes
 \mathcal C_W^{\,m(z\to z')}\psi.
 \label{eq:cpe-transporte-cilindrico-cayley}
\end{equation}
La proyectividad de \(\mu_k\), la ortogonalidad de descendientes y la
unitariedad de \(\mathcal C_W\) dan
\(\widetilde U_k^*\widetilde U_k=I\). La etiqueta completa \(z'\)
conserva cilindro, hoja, orientación, acarreo, ruta, frontera y
multiplicidad; por ello, historias distintas no se identifican al aplicar
\eqref{eq:cpe-transporte-cilindrico-cayley}.

\begin{theorem}[Entrelazamiento de Cayley bajo \(1\) vuelta nonádica completa]
\label{thm:cpe-entrelazamiento-cayley-tpk}
Si \(\Gamma_9\) es la holonomía nonádica completa, normalizada por
\(m(\Gamma_9)=1\), entonces
\begin{equation}
 \boxed{\widetilde U_{\Gamma_9,k}J_k
 =J_{k+9}\mathcal C_W.}
 \label{eq:cpe-entrelazamiento-exacto}
\end{equation}
Para todo átomo espectral \(\psi_\gamma\) de \(H_W\), con
\(H_W\psi_\gamma=\gamma\psi_\gamma\), se cumple
\begin{equation}
 \widetilde U_{\Gamma_9^n,k}J_k\psi_\gamma
 =\tau_\gamma^nJ_{k+9n}\psi_\gamma,
 \qquad
 \tau_\gamma=\frac{\gamma+i/2}{\gamma-i/2}\in S^1.
 \label{eq:cpe-orbita-espectral-ledger}
\end{equation}
En la publicación contractiva, la misma identidad toma la forma
\[
 \widetilde M_{\Gamma_9^n,k}J_k\psi_\gamma
 =\left(\frac59\right)^n\tau_\gamma^n
 J_{k+9n}\psi_\gamma.
\]
\end{theorem}

\begin{proof}
La vuelta \(\Gamma_9\) retorna a la fase visible pero incrementa en una
unidad la memoria, de modo que \(m(\Gamma_9)=1\). Al sumar todos los
descendientes de \(\Omega_k\), los pesos telescopan con la medida proyectiva
y producen \(\Omega_{k+9}\), mientras la fibra de Weil recibe exactamente
una potencia de \(\mathcal C_W\). Esto prueba
\(\widetilde U_{\Gamma_9,k}J_k=J_{k+9}\mathcal C_W\). La iteración usa la
aditividad \(m(\Gamma_9^n)=n\). La última fórmula sigue de
\(\widetilde M_{\Gamma_9,k}=(5/9)\widetilde U_{\Gamma_9,k}\).
\end{proof}

El teorema no identifica un cero con una fase aislada del calendario. Una
vuelta completa aplica una vez el carácter de Cayley; cada cero determina
la ley angular con la que una historia atraviesa todas las vueltas. En
particular, retorno de fase y retorno de estado permanecen distintos:
\[
 g(K+9)=g(K),
 \qquad B_{K+9}\ne B_K\quad\text{en general}.
\]

\section{Convergencia en norma de las matrices espectrales finitas}
\label{sec:cpe-convergencia-resolvente}

Sea
\[
 R_W=(H_W-iI)^{-1}.
\]
El espectro de \(H_W\) es discreto y \(|\gamma|\to\infty\), por lo que
\(R_W\) es compacto. Se toma una familia numerable
\(\{f_j\}_{j\ge1}\subset C_c^\infty(\mathbb R)\) con parámetros racionales,
densa en el núcleo de pruebas y elegida sin consultar alturas espectrales.
Después de cocientar el núcleo de la forma de Weil y ortonormalizar, se
obtienen subespacios
\[
 \mathcal V_N=\operatorname{span}\{[f_1],\ldots,[f_N]\},
 \qquad P_N:\mathcal H_W\to\mathcal V_N,
\]
con \(P_N\to I\) fuertemente. Defínase
\begin{equation}
 R_N=P_NR_WP_N.
 \label{eq:cpe-resolvente-finito}
\end{equation}

\begin{theorem}[Aproximación compacta sin introducir los ceros]
\label{thm:cpe-convergencia-norma}
La sucesión \(R_N\) converge en norma a \(R_W\):
\begin{equation}
 \boxed{\lim_{N\to\infty}\|R_W-R_N\|=0.}
 \label{eq:cpe-convergencia-norma}
\end{equation}
En consecuencia, cada autovalor no nulo
\(\lambda_\gamma=(\gamma-i)^{-1}\) de \(R_W\) es límite, con su
multiplicidad algebraica, de autovalores de las matrices finitas de
\(R_N\). Las alturas se recuperan mediante
\begin{equation}
 \gamma=i+\lambda_\gamma^{-1}.
 \label{eq:cpe-recuperacion-altura}
\end{equation}
\end{theorem}

\begin{proof}
Si \(K\) es compacto y \(P_N\to I\) fuertemente, la convergencia es uniforme
sobre el compacto \(K(B_1)\), donde \(B_1\) es la bola unidad. Por tanto,
\[
 \|(I-P_N)K\|\longrightarrow0.
\]
Aplicando el mismo argumento a \(K^*\) se obtiene
\(\|K(I-P_N)\|\to0\). Así,
\[
 \|K-P_NKP_N\|
 \leq\|(I-P_N)K\|+\|P_NK(I-P_N)\|\longrightarrow0.
\]
Tomando \(K=R_W\) se prueba \eqref{eq:cpe-convergencia-norma}. La
estabilidad de los autovalores aislados de operadores compactos bajo
perturbaciones en norma y el cálculo funcional de Riesz conservan su
multiplicidad. La identidad \eqref{eq:cpe-recuperacion-altura} es la inversa
de la aplicación espectral del resolvente.
\end{proof}

Las entradas de \(R_N\) se calculan desde la forma primo--gamma--polar:
\[
 \langle R_W[f_j],[f_k]\rangle_W
 =i\int_0^\infty e^{-t}
 \mathscr I_{\rm GW}(T_{-t}f_j,f_k)\,dt.
\]
Por ello, el teorema no utiliza una tabla de ceros para formar las matrices.

\subsection{Certificación a posteriori de intervalos y multiplicidades}

Sea \((\lambda_N,v_N)\) un autovector normalizado de \(R_N\) y sea
\[
 \varepsilon_N(v_N)
 :=\|(I-P_N)R_Wv_N\|.
\]
Como \(R_W\) es normal,
\begin{equation}
 \operatorname{dist}(\lambda_N,\sigma(R_W))
 \leq\varepsilon_N(v_N).
 \label{eq:cpe-residuo-inclusion}
\end{equation}
Los canales primo, gamma y polar proporcionan cotas intervalares separadas
para \(\varepsilon_N\); sus colas se suman antes de emitir el intervalo.
Además, si \(\Gamma\) es un contorno cerrado contenido en el resolvente de
\(R_N\) y
\begin{equation}
 \sup_{z\in\Gamma}
 \|(z-R_N)^{-1}\|\,\|R_W-R_N\|<1,
 \label{eq:cpe-criterio-riesz}
\end{equation}
entonces las proyecciones de Riesz de \(R_N\) y \(R_W\) delimitadas por
\(\Gamma\) poseen el mismo rango. Las ecuaciones
\eqref{eq:cpe-residuo-inclusion} y \eqref{eq:cpe-criterio-riesz} certifican,
respectivamente, inclusión y completitud de los átomos contenidos en cada
ventana. La aritmética intervalar endurece los decimales; no añade un mapa
matemático ausente.

\section{Traza acoplada del espectro de Weil y la graduación Moonshine}
\label{sec:cpe-traza-rh-moonshine}

La traza producto
\[
 \operatorname{Tr}_{\mathcal H_W}
 \bigl(e^{-aH_W^2}\mathcal C_W^n\bigr)
 \operatorname{Tr}_{V^\natural}
 \bigl(g e^{-\beta(L_0-1)}\bigr)
\]
superpone dos observables, pero no los acopla. La memoria del TPK y la
graduación de \(V^\natural\) permiten construir el acoplamiento sin postular
una acción del Monstruo sobre cifras ni sobre alturas.

Fijado el calibre excepcional \(\chi_{\rm exc}\), se orientan conjuntamente
el origen de memoria y el origen de la graduación. Sea
\(\mathcal H_{\rm mem}=\ell^2(\mathbb N_0)\), con
\(N_{\rm mem}|n\rangle=n|n\rangle\), y sea
\[
 V^\natural=\widehat\bigoplus_{r\ge0}V^\natural_r,
 \qquad L_0|_{V^\natural_r}=rI.
\]
Denotemos por \(P_r^\natural\) la proyección sobre
\(V^\natural_r\). El proyector de igualdad de grados es
\begin{equation}
 \Pi_{\chi_{\rm exc}}
 :=\sum_{n\ge0}|n\rangle\langle n|
 \otimes I_{\mathcal H_W}\otimes P_n^\natural.
 \label{eq:cpe-proyector-diagonal-grados}
\end{equation}
La suma converge fuertemente y define una proyección ortogonal. Sobre el
producto se define el carácter de memoria
\begin{equation}
 \mathcal C_W^{N_{\rm mem}}
 :=\sum_{n\ge0}|n\rangle\langle n|
 \otimes\mathcal C_W^n.
 \label{eq:cpe-cayley-memoria}
\end{equation}

\begin{definition}[Traza diagonal RH--Moonshine]
\label{def:cpe-traza-diagonal}
Para \(a,b,\beta>0\) y \(g\in\mathbb M\), se define
\begin{equation}
\begin{aligned}
 \Theta^{\Delta}_{g,\chi_{\rm exc}}(a,b,\beta)
 :={}&\operatorname{Tr}\!\left[
 \Pi_{\chi_{\rm exc}}
 \Bigl(
 e^{-bN_{\rm mem}}
 \mathcal C_W^{N_{\rm mem}}e^{-aH_W^2}
 \otimes g e^{-\beta(L_0-1)}
 \Bigr)
 \Pi_{\chi_{\rm exc}}
 \right].
 \label{eq:cpe-traza-diagonal-def}
\end{aligned}
\end{equation}
\end{definition}

\begin{theorem}[Existencia y expansión no factorizada]
\label{thm:cpe-traza-diagonal}
La traza \eqref{eq:cpe-traza-diagonal-def} es absolutamente convergente y
satisface
\begin{equation}
 \boxed{\Theta^{\Delta}_{g,\chi_{\rm exc}}(a,b,\beta)
 =\sum_{n\ge0}e^{-bn}e^{-\beta(n-1)}
 \operatorname{Tr}_{\mathcal H_W}
 \bigl(e^{-aH_W^2}\mathcal C_W^n\bigr)
 \operatorname{Tr}_{V^\natural_n}(g).}
 \label{eq:cpe-traza-diagonal-expansion}
\end{equation}
Después del reconocimiento espectral,
\begin{equation}
 \Theta^{\Delta}_{g,\chi_{\rm exc}}(a,b,\beta)
 =\sum_{n\ge0}\sum_\gamma
 m_\gamma e^{-a\gamma^2}
 e^{-bn-\beta(n-1)}\tau_\gamma^n
 \operatorname{Tr}_{V^\natural_n}(g).
 \label{eq:cpe-traza-diagonal-espectral}
\end{equation}
Esta expresión no es el producto de una suma espectral y una serie de
McKay--Thompson: el mismo índice \(n\) selecciona simultáneamente la potencia
del carácter de memoria y el grado monstruoso.
\end{theorem}

\begin{proof}
Para \(a>0\), la ley de conteo espectral implica
\(\sum_\gamma m_\gamma e^{-a\gamma^2}<\infty\). Cada espacio
\(V^\natural_n\) es finito-dimensional y los coeficientes de los caracteres
de Moonshine crecen subexponencialmente en \(n\). El factor
\(e^{-(b+\beta)n}\) asegura la convergencia absoluta de la serie conjunta.
La expansión se obtiene insertando las resoluciones espectrales de
\(N_{\rm mem}\), \(H_W\) y \(L_0\) y utilizando
\eqref{eq:cpe-proyector-diagonal-grados}. La presencia del proyector diagonal
elimina los términos con grados distintos y produce una sola suma en
\(n\), en vez de las dos sumas independientes de la traza producto.
\end{proof}

El acoplamiento depende del calibre \(\chi_{\rm exc}\), que fija origen,
orientación y alineamiento de la lectura excepcional. Una elección conjugada
produce una traza conjugada o una reindexación declarada; no modifica el
teorema de Riemann--Weil ni el entrelazamiento
\eqref{eq:cpe-entrelazamiento-exacto}.

